\scanpage{42}
% archive sheet number: 41a (upper slip)
% slanted
\margin{Le rôle de 3.1.13 est bien obscur. Ne peut-on le \unsure{vider|virer} avantageusement ?}
On a dans ce cas, posant $K^!=K'{}^!$,
\[
 K^{* !}\underline F=a\bigl(W\longmapsto\mathrm{Hom}_{V_W}(K(W),\underline F\mid V_W)\bigr),
\]
\[
 \underline H^n(RK^!L)_x\simeq\underset{\unsure{W(x)|W(\bar x)}}{\underrightarrow{\lim}}\,
 \mathrm{Ext}^n_{V_W}(K(W),L\mid V_W).
\]

% Lower sheet; a long pale oblique line crosses the writing from top to bottom.
\tikz[remember picture,overlay]{\coordinate (scanXLIIstart) at (1.3,0);}
% slanted
\margin{À quoi sert cette section 3.1.14, qu'y fait-on ?}
3.1.\deleted{\unsure{16|10}}14. Appliquons les considérations qui précèdent à un morphisme compactifiable de schémas $f:X\longrightarrow S$.
D'après les démonstrations de 3.1.\deleted{\unsure{11}}3 et 3.1.4, si $f=\bar f j$ est une compactification de $f$ et si les fibres de $f$ sont de dimension $\leq d$, le foncteur $Rf^!$ se construit par le procédé de 3.1.10, pour
\margin{3.1.10}
\[
 K(U)=\bar f_*\tau^{\unsure{\bullet|*}}_{\leq d}C^{*}_{\ell}\mathcal A_U=f_!^\bullet(\mathcal A_U).
\]
\margin{$\tau''_{\leq 2d}$ ?}
\inserted{où $A_U$ ---}
% slanted
\margin{dire qui est $A_U$}

Le corollaire 3.1.12 va permettre de donner d'autres procédés de calcul\unsure{/} du foncteur $Rf^!$. Supposons que $n\mathcal A=0$. On a alors ($\overline{\underline{\mathrm{XVII}}}$ 5.2\tikz[baseline]{\draw (0,-.07) -- (.06,.42); \draw (.02,.06) -- (.10,-.02);})
\margin{\tikz{\draw (0,-.07) -- (.06,.42); \draw (.02,.06) -- (.10,-.02);}}
\[
 Rf_!\mathcal A_U\simeq\mathcal A\overset{\mathbb L}{\otimes}Rf_!(\mathbb Z/n)_U.
\]
En fait, $f_!^\bullet((\mathbb Z/n)_U)$ est un complexe de faisceaux \uline{plats} : il suffit de le voir pour $n$\deleted{$=$}\inserted{de la forme} $\ell^k$ \deleted{primaire}, \deleted{\unsure{et}} la platitude \deleted{\unsure{degré par degré}} d'un faisceau $\underline F$ \deleted{\unsure{de gr}} de $\mathbb Z/\ell^k$-modules signifie que pour $i<k$, $\underline F/\ell\underline F\xrightarrow{\sim}\ell^i\underline F/\ell^{i+1}\underline F$, et ceci résulte pour les faisceaux $f_!^i(\mathbb Z/n)_U$ de l'exactitude de $f_!^i$ et de la platitude \deleted{\unsure{de $f_!^i$}} $(\mathbb Z/n)_U$.
\deleted{On a donc un quasi-isomorphisme} La flèche \inserted{canonique}
\[
 \deleted{f_!^\bullet(\mathbb Z/n)_U\xrightarrow{\sim}}
\]
\[
 f_!^\bullet(\mathbb Z/n)_U\otimes\mathcal A\xrightarrow{\sim}f_!^\bullet(\mathcal A_U)
\]
est donc un quasi-isomorphisme et on peut aussi bien décrire (3.1.12)
\tikz[remember picture,overlay]{\draw[gray,thin] (scanXLIIstart) -- (0,0);}
% archive sheet number: 41
